\documentclass[11pt]{article}

\usepackage[margin=1in]{geometry}
\usepackage{amsmath}
\usepackage{amssymb}
\usepackage{booktabs}
\usepackage{graphicx}
\usepackage{hyperref}
\usepackage{microtype}
\usepackage{xcolor}

\hypersetup{
  colorlinks=true,
  linkcolor=blue,
  citecolor=blue,
  urlcolor=blue
}

\title{A Tiny Interpretable Codec for Analogy Structure in Qwen Hidden States}
\author{Larry Richards}
\date{\today}

\begin{document}

\maketitle

\begin{abstract}
This report describes a small experiment in extracting a human-readable semantic
coordinate system from the hidden states of a frozen language model. We take
hidden activations from Qwen 2.5 for four words---\emph{king}, \emph{queen},
\emph{man}, and \emph{woman}---and train a linear encoder that maps those
high-dimensional activations into a two-dimensional latent space. In the learned
space, the classic analogy relation
\[
  \mathrm{king} - \mathrm{man} + \mathrm{woman} \approx \mathrm{queen}
\]
holds with cosine similarity around 0.9996. The four word classes are also
perfectly separated by a nearest-mean classifier in the two-dimensional space,
and the learned gender and royalty directions are nearly orthogonal. The point
of the experiment is not to claim a universal semantic basis inside Qwen. It is
to show, in a controlled setting, that useful semantic structure can be pulled
out of model activations and compressed into a very small, inspectable interface.
\end{abstract}

\section{Introduction}

Large language models represent words and contexts as high-dimensional internal
states. These states are powerful, but they are not naturally easy to inspect.
The broad question behind this experiment is simple: can we build a small
interface, or \emph{codec}, that exposes a meaningful slice of this internal
structure in a form a human can look at?

This experiment focuses on one deliberately tiny semantic domain: the familiar
analogy
\[
  \mathrm{king} - \mathrm{man} + \mathrm{woman} \approx \mathrm{queen}.
\] \cite{mikolov2013analogies}.
Rather than testing this relation in an embedding table, we extract hidden
states from a frozen Qwen 2.5 instruction model \cite{qwen2024} while the
target words appear in simple sentence contexts. We then train a
two-dimensional linear projection that organizes those hidden states into an
interpretable plane.

The result is a compact proof of concept. A two-dimensional learned space is
enough to preserve the analogy, separate the four classes, and expose two
nearly orthogonal semantic directions corresponding roughly to gender and
royalty. This is a small result, but it is useful because it gives the larger
``latent bus'' idea something concrete to stand on: model internals can be
wrapped with small, testable semantic interfaces.

\section{Experiment Overview}

The experiment has three stages.

\begin{enumerate}
  \item Extract hidden-state representations for the target words from a frozen
  Qwen model.
  \item Train a linear codec from the model's hidden dimension into a
  two-dimensional latent space.
  \item Validate the learned space using analogy, classification, and axis
  geometry checks.
\end{enumerate}

The target vocabulary is intentionally small:
\[
  \mathcal{W} = \{\mathrm{king}, \mathrm{queen}, \mathrm{man}, \mathrm{woman}\}.
\]
For each word, the extraction script generates 100 short sentence contexts from
simple templates such as:

\begin{quote}
I saw the king today. \\
The king was there. \\
Look at the king. \\
Where is the king?
\end{quote}

The target word is located using tokenizer offset mappings, and the hidden state
for the final token of the word is extracted from a selected model layer. In
the original run, the selected hidden layer used an offset of 10 from the final
layer. This gives 400 total hidden-state examples: 100 for each of the four
words.

\section{Codec Model}

Let \(x \in \mathbb{R}^d\) be a Qwen hidden-state vector for a target word in
context. The codec learns a linear encoder
\[
  z = x W_e + b_e,
\]
where \(z \in \mathbb{R}^2\). The purpose of the encoder is not merely to reduce
dimension. It is trained to produce a two-dimensional space where the analogy
geometry and class separation are both explicit. In this respect it is close in
spirit to linear probing, but with an explicit geometric target rather than
only a diagnostic classifier \cite{alain2016, belinkov2022}.

The training objective combines several terms:

\[
  \mathcal{L}
  =
  \mathcal{L}_{cls}
  + 3.0 \mathcal{L}_{analogy}
  + \mathcal{L}_{perp}
  + 0.1 \mathcal{L}_{orth}
  + 0.5 \mathcal{L}_{sep}.
\]

Here, \(\mathcal{L}_{cls}\) is a classification loss in the two-dimensional
space. The analogy loss encourages
\[
  \mu_{\mathrm{king}} - \mu_{\mathrm{man}} + \mu_{\mathrm{woman}}
  \approx
  \mu_{\mathrm{queen}},
\]
where \(\mu_w\) is the class mean for word \(w\) after projection. The
perpendicularity term encourages the gender and royalty directions to be close
to orthogonal. A light encoder regularizer keeps the projection from becoming
pathological, and a separation term discourages the four class means from
collapsing together.

This is deliberately simple: no nonlinear encoder, no complicated probing
architecture, and no model fine-tuning. The language model is frozen; only the
small projection and a lightweight classifier head are trained.

\section{Validation}

The experiment uses three main validation gates.

\subsection{Analogy Gate}

The first gate checks whether the analogy holds in the learned two-dimensional
space:
\[
  \cos(
    \mu_{\mathrm{king}} - \mu_{\mathrm{man}} + \mu_{\mathrm{woman}},
    \mu_{\mathrm{queen}}
  ) \geq 0.95.
\]

In the saved run, this value is approximately 0.9996. This is the strongest
single sign that the codec is organizing the four concepts into the intended
geometry.

\subsection{Classification Gate}

The second gate asks whether individual projected examples can be classified
correctly in the two-dimensional space. The validation code uses a simple
nearest-mean classifier. Each held-out projected point is assigned to whichever
class mean is closest in Euclidean distance:
\[
  \hat{y}(z) = \arg\min_{w \in \mathcal{W}} \|z - \mu_w\|_2.
\]

The data points used here are not synthetic perturbations in the two-dimensional
space. They are real hidden-state activations extracted from different sentence
contexts, then projected through the learned codec. The original validation uses
a 70/30 split per class for the classification report, producing 30 held-out
examples per word and 120 held-out examples total.

\subsection{Axis Geometry Gate}

The third gate checks whether the two main semantic directions are nearly
orthogonal. The gender direction is defined as
\[
  g = \mu_{\mathrm{woman}} - \mu_{\mathrm{man}},
\]
and the royalty direction is defined as
\[
  r = \mu_{\mathrm{king}} - \mu_{\mathrm{man}}.
\]

The angle between these two directions is approximately \(89.7^\circ\), which
is very close to the intended right angle.

\section{Results}

Table~\ref{tab:main-results} summarizes the main validation results.

\begin{table}[h]
\centering
\begin{tabular}{llll}
\toprule
Gate & Metric & Result & Target \\
\midrule
Analogy & Cosine similarity & 0.9996 & \(\geq 0.95\) \\
Classification & Held-out accuracy & 100.0\% & \(\geq 98\%\) \\
Axis geometry & Gender/royalty angle & \(89.7^\circ\) & near \(90^\circ\) \\
\bottomrule
\end{tabular}
\caption{Main validation gates for the trained two-dimensional codec.}
\label{tab:main-results}
\end{table}

The held-out confusion matrix is perfectly diagonal:

\begin{table}[h]
\centering
\begin{tabular}{lrrrr}
\toprule
True class & king & queen & man & woman \\
\midrule
king  & 30 & 0  & 0  & 0 \\
queen & 0  & 30 & 0  & 0 \\
man   & 0  & 0  & 30 & 0 \\
woman & 0  & 0  & 0  & 30 \\
\bottomrule
\end{tabular}
\caption{Nearest-mean classification in the learned two-dimensional space.}
\label{tab:confusion}
\end{table}

The important thing about this result is not that the classifier is complex. It
is the opposite: classification is almost embarrassingly simple once the codec
has organized the representations. The model activations land in a plane where
the four concepts are cleanly separated, and the analogy is visible as geometry.

\subsection{Loss Ablation}

The full objective deliberately asks for both class separation and analogy
geometry. To check how much of the result is imposed by the loss, we reran the
codec training under five loss configurations while holding the data, seed,
optimizer, learning rate, and number of steps fixed. Table~\ref{tab:loss-ablation}
shows the result.

\begin{table}[h]
\centering
\begin{tabular}{lrrr}
\toprule
Training objective & Analogy cosine & Accuracy & Axis angle \\
\midrule
Full loss & 1.0000 & 100.0\% & \(90.0^\circ\) \\
Drop \(\mathcal{L}_{analogy}\) & 0.9951 & 100.0\% & \(90.0^\circ\) \\
Drop \(\mathcal{L}_{perp}\) & 0.9998 & 100.0\% & \(98.5^\circ\) \\
Drop \(\mathcal{L}_{analogy}\) and \(\mathcal{L}_{perp}\) & 0.9971 & 100.0\% & \(96.4^\circ\) \\
Classification only & 0.9840 & 100.0\% & \(105.8^\circ\) \\
\bottomrule
\end{tabular}
\caption{Loss ablation with fixed seed, data, optimizer, learning rate, and
3000 training steps. The ablation uses the saved extracted activations and
re-trains the small codec for each configuration.}
\label{tab:loss-ablation}
\end{table}

The interesting result is the last row. Even with classification loss alone,
the learned two-dimensional space still gives an analogy cosine of 0.9840. The
full objective sharpens the geometry and makes the axes nearly orthogonal, but
the analogy relation is not created entirely by the analogy term. Some of the
structure is already available in the activations and can be recovered by a
very small supervised interface.

\section{What This Shows}

The experiment supports three modest claims.

First, the Qwen hidden states contain enough information to distinguish these
four word classes across simple contexts. This is expected, but the clean
two-dimensional separation makes it directly visible.

Second, the analogy structure can be made explicit with a small supervised
projection. The relation does not have to be recovered by inspecting the full
hidden state directly; it can be exposed through a tiny learned interface.

Third, the learned space is not just a classifier trick. The geometry has
interpretable structure: the analogy closes, and the two hand-named semantic
directions are close to orthogonal.

This is why the experiment is interesting. It suggests that at least some
semantic circuits inside a model can be treated as something like an API: a
small codec translates opaque activation vectors into coordinates that can be
tested, plotted, and reasoned about. This framing is compatible with the broader
linear representation hypothesis: useful semantic variables may sometimes be
represented in directions or low-dimensional subspaces that simple maps can
recover \cite{park2023}.

\section{Limitations}

This is a proof of concept, not a broad interpretability result.

The vocabulary is tiny. Four words are enough to test the idea, but not enough
to claim that the same two axes generalize across language. A stronger version
would include many royal terms, many gendered terms, neutral controls, and
examples across languages and syntactic settings.

The classification gate is also not a strict generalization benchmark. The
codec is trained using the full extracted dataset, and the held-out split is
used afterward to check separability in the learned space. This is acceptable
for a first demonstration, but future work should use a cleaner train,
validation, and test protocol. More generally, probe results need controls:
a sufficiently expressive probe can sometimes learn the task rather than reveal
what was already explicit in the representation \cite{hewitt2019}.

The sentence templates are simple. That is useful because it keeps the
experiment controlled, but it also means the result may depend on a narrow
distribution of contexts. More varied natural sentences would make the test
more realistic.

Finally, the axis names should be read carefully. Calling one direction
``gender'' and another ``royalty'' is a useful interpretation of this particular
four-word geometry. It is not evidence, by itself, that the model has a single
global gender direction or royalty direction in the larger hidden space.
Superposition also gives a reason to expect useful features to be entangled
rather than cleanly isolated in general \cite{elhage2022}.

\section{Next Steps}

The most natural next step is to scale the vocabulary while keeping the codec
small. For example, one could include terms such as \emph{prince},
\emph{princess}, \emph{duke}, \emph{duchess}, \emph{boy}, \emph{girl},
\emph{father}, and \emph{mother}, plus neutral controls. If the same two axes
remain useful, the result becomes much more compelling.

Another direction is to test cross-context generalization. The codec could be
trained on simple templates and evaluated on natural text, or trained on English
and tested on translated equivalents. This would help distinguish a real
semantic coordinate system from a template-specific artifact.

A third direction is intervention. If a point can be moved along a learned axis
in the two-dimensional space, can the corresponding hidden-state intervention
produce predictable changes in model behavior? That would move the codec from
being a visualization tool toward being a steering interface. This would connect
the codec to causal abstraction and model-editing work, where the goal is not
only to read a representation but to test whether changing it changes behavior
\cite{geiger2021, meng2022, nanda2023}.

Another useful comparison point is unsupervised dictionary learning, which aims
to recover sparse interpretable features from model activations at larger scale
\cite{templeton2024}. The present codec is much smaller and supervised, but it
shares the same practical ambition: turning opaque activation spaces into
interfaces with named, testable coordinates.

\section{Conclusion}

This experiment shows that a very small linear codec can extract a clean,
interpretable analogy space from frozen Qwen hidden states. In that space, the
relation \(\mathrm{king} - \mathrm{man} + \mathrm{woman} \approx
\mathrm{queen}\) holds almost exactly, the four classes are perfectly separable
under a simple nearest-mean classifier, and the main semantic directions are
nearly orthogonal.

The result is small but suggestive. It points toward a practical style of
mechanistic work where model internals are not only probed, but wrapped in
compact interfaces with explicit validation gates. That is the larger promise:
not just looking at activations, but building little semantic instruments that
make them easier to test, compare, and eventually use.

\begin{thebibliography}{99}

\bibitem{mikolov2013}
T.~Mikolov, K.~Chen, G.~Corrado, and J.~Dean.
\newblock Efficient estimation of word representations in vector space.
\newblock \emph{arXiv preprint arXiv:1301.3781}, 2013.

\bibitem{mikolov2013analogies}
T.~Mikolov, W.~Yih, and G.~Zweig.
\newblock Linguistic regularities in continuous space word representations.
\newblock In \emph{Proceedings of NAACL-HLT}, pages 746--751, 2013.

\bibitem{alain2016}
G.~Alain and Y.~Bengio.
\newblock Understanding intermediate layers using linear classifier probes.
\newblock \emph{arXiv preprint arXiv:1610.01644}, 2016.

\bibitem{hewitt2019}
J.~Hewitt and P.~Liang.
\newblock Designing and interpreting probes with control tasks.
\newblock In \emph{Proceedings of EMNLP-IJCNLP}, pages 2733--2743, 2019.

\bibitem{belinkov2022}
Y.~Belinkov.
\newblock Probing classifiers: Promises, shortcomings, and advances.
\newblock \emph{Computational Linguistics}, 48(1):207--219, 2022.

\bibitem{park2023}
K.~Park, Y.~J. Choe, and V.~Veitch.
\newblock The linear representation hypothesis and the geometry of large language models.
\newblock \emph{arXiv preprint arXiv:2311.03658}, 2023.

\bibitem{elhage2022}
N.~Elhage, T.~Hume, C.~Olsson, N.~Schiefer, T.~Henighan, S.~Kravec,
  Z.~Hatfield-Dodds, R.~Lasenby, D.~Drain, C.~Chen, R.~Grosse, S.~McCandlish,
  J.~Kaplan, D.~Amodei, M.~Wattenberg, and C.~Olah.
\newblock Toy models of superposition.
\newblock \emph{Transformer Circuits Thread}, 2022.

\bibitem{geiger2021}
A.~Geiger, H.~Lu, T.~Icard, and C.~Potts.
\newblock Causal abstractions of neural networks.
\newblock In \emph{Advances in Neural Information Processing Systems}, volume~34, pages 9574--9586, 2021.

\bibitem{meng2022}
K.~Meng, D.~Bau, A.~Andonian, and Y.~Belinkov.
\newblock Locating and editing factual associations in GPT.
\newblock In \emph{Advances in Neural Information Processing Systems}, volume~35, pages 17359--17372, 2022.

\bibitem{nanda2023}
N.~Nanda, A.~Lee, and M.~Wattenberg.
\newblock Emergent linear representations in world models of self-supervised sequence models.
\newblock \emph{arXiv preprint arXiv:2309.00941}, 2023.

\bibitem{templeton2024}
A.~Templeton, T.~Conerly, J.~Marcus, J.~Lindsey, T.~Bricken, B.~Chen,
  A.~Pearce, C.~Citro, E.~Ameisen, A.~Jermyn, T.~Henighan, S.~Kravec,
  Z.~Hatfield-Dodds, R.~Lasenby, D.~Hume, J.~Bills, B.~Lee, N.~Schiefer,
  J.~Batson, M.~Wattenberg, and C.~Olah.
\newblock Scaling monosemanticity: Extracting interpretable features from
  Claude 3 Sonnet.
\newblock \emph{Transformer Circuits Thread}, 2024.

\bibitem{qwen2024}
A.~Yang, B.~Yang, B.~Hui, et~al.
\newblock Qwen2 technical report.
\newblock \emph{arXiv preprint arXiv:2407.10671}, 2024.

\end{thebibliography}

\end{document}
